% ECONOMICS_PROBLEM: {"id": "D72-476", "title": "Causes and economic consequences of democratic erosion", "jel_code": "D72", "source_id": "OP-0121"}
% !TeX program = xelatex
\documentclass[11pt,a4paper]{article}
\usepackage[margin=23mm]{geometry}
\usepackage{fontspec}
\setmainfont{Latin Modern Roman}
\usepackage{amsmath,amssymb,mathtools}
\usepackage{xcolor,enumitem,booktabs,longtable,array,xurl,hyperref}
\definecolor{muted}{HTML}{53636C}
\hypersetup{colorlinks=true,urlcolor=blue,linkcolor=blue}
\setlength{\parindent}{0pt}
\setlength{\parskip}{5pt}
\newcommand{\R}{\mathbb R}
\newcommand{\E}{\mathbb E}
\newcommand{\N}{\mathbb N}
\newcommand{\ind}{\mathbf 1}
\newcommand{\Var}{\operatorname{Var}}
\newcommand{\Cov}{\operatorname{Cov}}
\newcommand{\supp}{\operatorname{supp}}
\DeclareMathOperator*{\argmin}{arg\,min}
\DeclareMathOperator*{\argmax}{arg\,max}
\newcommand{\entrymeta}[3]{{\sffamily\small\color{muted}\textbf{\texttt{#1}}\quad|\quad #2\par #3\par}}
\newcommand{\Problemref}[1]{\texttt{#1}}
\newcommand{\JELref}[2]{\texttt{#2} (JEL \texttt{#1})}
\begin{document}
\section*{Causes and economic consequences of democratic erosion}
\textbf{Problem ID:} \texttt{D72-476}\quad \textbf{JEL:} \texttt{D72}\par
% BEGIN ECONOMICS STATEMENT
\label{problem:OP-0121}
\label{atlas:prob:Q1}
{\sffamily\small\color{muted}\textbf{Q1} | Political economy | Editorial research problem family\par}
\subsection*{Definitions and assumptions}
Fix countries $c$ in a declared sampling frame and dates $t=0,\ldots,T$. Let $D_{ct}\in\{0,1\}$ be a prespecified democracy classification, $J_{ct}$ a vector of executive constraints, and $Y_{ct}>0$ real output per capita. An intervention $a$ specifies a feasible economic-shock process or institutional protection, its implementation date, and its effects on the government's available actions. Let $(D_{ct}(a),J_{ct}(a),Y_{ct}(a))$ be the potential history, allowing interference across countries through a stated joint policy vector. A structural model $m\in\mathcal M$ specifies citizen and elite decisions, state transitions, an observation process, and an equilibrium selection rule. Require finite expectations of $|\log Y_{ct}(a)|$.

For a democratic baseline date $t$, define the first subsequent loss of democracy by $\tau_c(a)=\inf\{s>t:D_{cs}(a)=0\}$, with $\inf\varnothing=\infty$. This operational definition must be supplemented if erosion within a still-democratic regime is the outcome of interest. The compared interventions share the same realized history through baseline date $t$ and begin strictly afterward; $D_{ct}=1$ is therefore a common pre-intervention eligibility condition.
\subsection*{Formal research question}
For a fixed horizon $h\leq T-t$ and interventions $a,a_0$, identify, or sharply bound over observationally compatible $m\in\mathcal M$,
\[
\begin{aligned}
\theta_h&=\mathbb E\!\left[\mathbf1\{\tau_c(a)\leq t+h\}-\mathbf1\{\tau_c(a_0)\leq t+h\}\mid D_{ct}=1\right],\\
\gamma_h&=\mathbb E[\log Y_{c,t+h}(a)-\log Y_{c,t+h}(a_0)\mid D_{ct}=1].
\end{aligned}
\]
Determine which restrictions on anticipation, political responses, common shocks and transition dynamics allow these quantities to be recovered from the declared data. Compare specified institutional interventions rather than treating the regime classification itself as an automatically well-defined treatment.
\subsection*{Interpretation and scope}
This is an editorial empirical and dynamic-equilibrium problem family. Democratization and democratic erosion need not have opposite effects of equal magnitude. A conditional transition probability in observed data is not the causal hazard under a reform. Persistent erosion requires a declared duration or institutional threshold; it cannot be inferred from one election result. Identification assumptions and regime measurement should be varied explicitly, with uncertainty propagated into both transition and output estimates.
\subsection*{Source audit and status}
The atlas's Acemoglu et al.\ (2019) and Guriev--Treisman (2019) are identifiable, but the former concerns democratization and the latter authoritarian information control. Neither explicitly poses the displayed equations. ``Democratic-transition literature'' is not an identifiable third citation; [S3] is a stated editorial replacement. This broad question is retained as a research agenda, with no claim that all transition models or settings are unresolved.
\subsection*{References}
{\footnotesize\raggedright
\begin{enumerate}[label={\textbf{[S\arabic*]}},leftmargin=3em]
\item Daron Acemoglu, Suresh Naidu, Pascual Restrepo, and James A. Robinson (2019). \emph{Democracy Does Cause Growth}. Journal of Political Economy 127(1), 47--100.\par
\textit{Locator and source role:} Abstract and dynamic-panel identification design; estimates democratization effects, not a general law of democratic erosion.\par
\url{https://www.journals.uchicago.edu/doi/10.1086/700936}
\item Sergei Guriev and Daniel Treisman (2019). \emph{Informational Autocrats}. Journal of Economic Perspectives 33(4), 100--127.\par
\textit{Locator and source role:} Abstract and discussion of information control; mechanisms of authoritarian rule, not causal identification of every regime transition.\par
\url{https://www.aeaweb.org/articles?id=10.1257/jep.33.4.100}
\item Daron Acemoglu and James A. Robinson (2001). \emph{A Theory of Political Transitions}. American Economic Review 91(4), 938--963.\par
\textit{Locator and source role:} Abstract and political-transition model; explicit editorial replacement for the unspecified transition literature.\par
\url{https://www.aeaweb.org/articles?id=10.1257/aer.91.4.938}
\end{enumerate}
}

\subsection*{Common conventions: Source mathematical conventions}

Unless an entry states otherwise, agent and firm sets are finite. State and action spaces are measurable, strategies and outcome maps are measurable, probability laws are specified, and expectations exist for the targets under discussion. Infinite-horizon discount factors lie in $(0,1)$ and discounted objectives are finite. Norms, tolerances and complexity input sizes must be specified for computational comparisons. Constants or primitives that appear locally are inputs, not empirical facts asserted by this catalogue.

\textbf{Identification.} A statistical model $m\in\mathcal M$ implies an observable law $P_m$ and target $\theta(m)$. At law $P$, its identified set is
\[
\Theta(P;\mathcal M)=\{\theta(m):m\in\mathcal M,\ P_m=P\}.
\]
Point identification holds when this set is a singleton, or equivalently when $P_{m_1}=P_{m_2}$ implies $\theta(m_1)=\theta(m_2)$ for all compatible models. Sharp bounds mean bounds attained, or approached when the boundary is not attained, by compatible models. Writing this set is a definition, not a solution: the substantive task is to establish informative restrictions and derive or estimate the set. An unrestricted-confounding model may give only trivial bounds.

\textbf{Inference.} Identification, estimability and valid uncertainty quantification are separate questions. Fix $\alpha\in(0,1)$. A confidence procedure $C_n$ over a declared law class $\mathcal P$ has asymptotic uniform coverage at level $1-\alpha$ when
\[
\liminf_{n\to\infty}\inf_{P\in\mathcal P}\Pr_P\{\theta(P)\in C_n\}\geq1-\alpha.
\]
For set-identified targets the coverage event must be defined separately, for example coverage of the full identified set. If an entry uses identification-indexed models instead of uniquely indexed laws, the same coverage convention is applied over those models. Pointwise limits do not establish uniform validity, and asymptotic validity does not guarantee finite-sample precision. Sample sizes, dependence structures and weak-identification sequences are part of each application.

\textbf{Counterfactuals.} $Y_i(a)$ is a potential outcome under a specified intervention, including its timing, implementation and versions. It is not $Y_i$ conditional on voluntarily choosing $a$. A full intervention vector permits interference; shorthand scalar treatment notation does not silently impose no interference. Assignment, selection, attrition, anticipation and measurement must be specified. A claim about an instrument's local effect is not automatically a population policy effect.

\textbf{Economic equilibrium.} A counterfactual model includes best responses, constraints, feasibility, market clearing, state transitions and expectations consistent with the stated information. When several equilibria exist, either a declared selector is an input or the target is set-valued. Comparisons across policy regimes must use the stated selection convention. Reduced-form relationships are not presumed invariant after an equilibrium-changing intervention.

\textbf{Optimization and welfare.} Optimization is over the stated feasible set. If attainment is not established, $\sup$ or $\inf$ and $\varepsilon$-optimal choices replace an asserted maximizer or minimizer. For example, with value $V=\sup_{a\in\mathcal A}W(a)<\infty$, an $\varepsilon$-optimal set is $\{a\in\mathcal A:W(a)\geq V-\varepsilon\}$ for $\varepsilon>0$. Benefits, costs, risk and health or environmental outcomes can be aggregated only using declared units and conversion weights. Interpersonal utility weights, the treatment of fiscal transfers and foreign profits, discounting, and the behavioral welfare criterion are explicit inputs. A comparative-static derivative additionally requires existence and appropriate differentiability of the selected counterfactual.

\textbf{Sources.} Local labels [S1], [S2], and so forth restart in every question. Locators identify the relevant abstract, model, section or result at the level actually checked; a bibliographic or abstract-level verification is not represented as inspection of an entire proof. Working papers are distinguished from later publications and changing versions. Source summaries are paraphrased. All URL checks and editorial formulations refer to the audit date stated above.
% END ECONOMICS STATEMENT
\end{document}
